% Source: https://artofproblemsolving.com/community/c3h2003706_questions
% Stable identity: aops-topic-2003706-q03
Put $A=(0,0)$, $B=(L,0)$, $C=(L,45)$, $D=(0,45)$ with $L=AB$. A circle of radius $10$ inside the rectangle tangent to $AB$ and $BC$ sits in corner $B$ and has centre $O_1=(L-10,10)$. The circle of radius $15$ in corner $D$ has centre $O_2=(15,30)$.

Two circles lying side by side inside the rectangle are tangent externally, so $O_1O_2=10+15=25$ (internal tangency would need $O_1O_2=5$, which is impossible because the vertical distance alone is $20$). Hence
\[
(L-25)^2+20^2=25^2\quad\Longrightarrow\quad (L-25)^2=225\quad\Longrightarrow\quad L=40\ \text{or}\ L=10 .
\]

\begin{asy}
usepackage("amsmath");
size(5cm);
defaultpen(fontsize(10pt));
real L=40;pair A=(0,0),B=(L,0),C=(L,45),D=(0,45);
pair P=(L-10,10),Q=(15,30);
draw(A--B--C--D--cycle,linewidth(1));
draw(circle(P,10));draw(circle(Q,15));
draw(P--Q,dashed);draw(P--(P.x,Q.y)--Q,dotted);
dot(P);dot(Q);dot(P+10*unit(Q-P));
label("$A$",A,SW);label("$B$",B,SE);label("$C$",C,NE);label("$D$",D,NW);
label("$O_1$",P,SE);label("$O_2$",Q,W);
label("$25$",(P+Q)/2,NE);label("$20$",(P.x,(P.y+Q.y)/2),E);label("$15$",((P.x+Q.x)/2,Q.y),N);
label("$45$",(L,22.5),E);
label("\mbox{}",(min(currentpicture,user=true)+max(currentpicture,user=true))/2,black); // keeps the text after the figure black when LaTeX typesets the labels inline
\end{asy}

If $L=10$, the circle of radius $15$ would not fit inside a rectangle of width $10$. With $L=40$ both circles lie inside: $O_1=(30,10)$ and $O_2=(15,30)$, and they touch. Therefore $AB=\boxed{40}$.
